\chapter{An empirical evaluation of log-likelihood}

% Tables extracted verbatim from arXiv:2409.06364v2 source.
% What happens to diffusion model likelihood when your model is conditional?
% Mattias Cross and Anton Ragni
% https://arxiv.org/src/2409.06364v2
%
% This is a LaTeX body fragment: use \input inside your document.
% All section/subsection headings are retained, including empty sections,
% so original section numbers (4.1, 4.2, Appendix B) are preserved.
% Tables 1--7 retain original data, captions, labels and source comments.
% Entirely commented-out draft tables are excluded.
%
% Required preamble:
% \usepackage{booktabs,multirow,siunitx,amssymb,pifont}
% \providecommand{\xmark}{\ding{55}}
%
% Source quirks retained: Table 4 has its label before its caption and
% surplus column declarations; Table 7 includes three reference columns
% that extend beyond the right edge of the published PDF.

\section{Data}
\paragraph{Speech}
The datasets considered are the validation sets of LJSpeech \citep{LJSpeechDataset} and TED-LIUM \citep{rousseauTEDLIUMAutomaticSpeech2012}. The LJSpeech split contains around 1 hour of audiobooks read by a US-English female speaker. The TED-LIUM split contains around 3 hours of male speakers and 1 hour of female speaker TED conference audio.
%
To observe how likelihood changes during the denoising diffusion process, the LJSpeech set is generated and likelihood is measured at various intervals of $t$. Likelihood is expected to increase as $t$ decreases. Likelihood is measured in \ac{BPD}, $-\log p(\x)\log_2\exp \cdot (\Pi_{i} \mathbf{d}_i)^{-1}$, where $d$ is the shape of $\x$. Lower BPD means higher likelihood.

\paragraph{Image}
Datasets considered are PACS \citep{liDeeperBroaderArtier2017}, a domain generalisation dataset that includes images of 9 classes in 4 domains (photo, art, cartoon, sketch) and CLEVR \citep{johnsonCLEVRDiagnosticDataset2017}, a diagnostic dataset containing synthetic images of 3 3D objects of 8 colours. The images used are generated from \cite{lewisDoesCLIPBind2023}, including a ``red sphere'' example. A subset of CLEVR consisting of single objects is used. For both datasets, each class-domain/shape-colour pair is repeated 20 times. Images are resized to $512^2$/$1024^2$.
The likelihood of SDXL is reported, with \ac{CFG} used to generate PACS. Image quality and compatibility are measured using \ac{FID} and CLIP similarity, respectively (Table~\ref{tab:sdxl_over_t}).

\section{Evaluation}
\paragraph{Speech}
\Ac{MCD} is a measure of how different two Mel-cepstra (audio features) are from each other \citep{kominekSynthesizerVoiceQuality2008}.\footnote{MCD is calculated with the method presented in ESPnet \citep{gaoEUROESPnetUnsupervised2023}} The root mean squared error of the log fundamental frequencies is denoted as ``LogF0'' and measures pitch/intonation accuracy.\footnote{For reference, Tacotron 2 and Fastspeech 2 achieve 0.26 and 0.24 LogF0 respectively \citep{renFastSpeechFastHighQuality2022}.} \Ac{WER} is the ratio of transcription errors to the number of words; intelligibility is measured by calculating the \ac{WER} of the transcript generated by the \ac{ASR} foundation model Whisper \citep{radfordRobustSpeechRecognition2022}. Speaker similarity with automatic speaker verification (ASV) is measured by comparing the cosine similarity between embeddings from the speaker verification foundation model Titanet \citep{koluguriTitaNetNeuralModel2021}.

\paragraph{Image}
FID measures the distance between real and synthetic image distributions. CLIP measures the compatibility between prompts and images. These metrics can be related to MCD and WER in the TTS experiment.


\section{Experiments}

\subsection{Text-to-speech}
\label{sec:Grad-TTS encoder and decoder output}

\begin{table}[htbp]
    \centering
    \caption{Quantitative assessment of Grad-TTS from time 1 to 0.}%Various metrics as time progresses from 1 to 0. These are likelihood (BPD), intonation error (LogF0), spectral distortion (MCD), intelligibility (WER), and speaker similarity (ASV). These results show that likelihood is correlated with smoother spectrograms and speaker accurate speech, but the diffusion generative process decreases intelligibility. There is a negligible change in intonation error.}
    \begin{tabular}{l||l|rrrrrrrrr}
        \toprule
        \multirow{2}*{Metric} & \multirow{2}*{Ref} & \multicolumn{9}{c}{Time}\\
         &  & 1.00 & 0.88 & 0.75 & 0.62 & 0.50 & 0.38 & 0.25 & 0.12 & 0.00 \\
        \midrule
        BPD $\downarrow$ & -0.06 & 2.79 & 2.79 & 2.79 & 2.79 & 2.77 & 2.72 & 2.55 & 1.35 & -0.08 \\
        LogF0 $\downarrow$ & 0.00 & 0.28 & 0.28 & 0.28 & 0.28 & 0.28 & 0.27 & 0.27 & 0.27 & 0.27 \\
        MCD $\downarrow$ & 0.00 & 6.52 & 6.52 & 6.52 & 6.50 & 6.47 & 6.43 & 6.39 & 6.34 & 6.31 \\
        WER $\downarrow$ & 0.03 & 0.08 & 0.08 & 0.08 & 0.09 & 0.08 & 0.10 & 0.13 & 0.16 & 0.19 \\
        ASV $\uparrow$ & 1.00 & 0.74 & 0.74 & 0.73 & 0.73 & 0.74 & 0.75 & 0.78 & 0.84 & 0.86 \\
        \bottomrule
    \end{tabular}
\label{tab:gtts_over_t}
\end{table}
To observe how likelihood changes during the denoising diffusion process, the LJSpeech set is generated and likelihood is measured at various intervals of $t$. Likelihood is expected to increase as $t$ decreases. Likelihood is measured in \ac{BPD}, $-\log p(\x)\log_2\exp \cdot (\Pi_{i} \mathbf{d}_i)^{-1}$, where $d$ is the shape of $\x$. Lower BPD means higher likelihood.
%
Table \ref{tab:gtts_over_t} shows that although likelihood increases as $\x_T$ is denoised into data $\x_0$, the effect on other metrics is surprising. 
Likelihood is correlated with ASV, marginally with MCD, but not LogF0. Intelligibility decreases as likelihood improves. The results suggest that composition between the source text and hypothesis utterance is modelled by the encoder $\enc(\cdot)$, and the diffusion decoder does not model the exact likelihood between text and speech $p_0(\x | \vy)$. This is verified through an \ac{ASR} decoder rescoring experiment on LJSpeech. An ASR pipeline generates a list of hypothesis transcripts for a given utterance, which are scored by a weighted combination of an \ac{AM} probability and a language model probability to find the best transcription. This pipeline serves as a diagnostic tool for how well Grad-TTS likelihoods can measure speech/text compatibility. The \ac{AM} probabilities are replaced with Grad-TTS likelihoods, and standard hypothesis rescoring techniques are applied \citep{kahnFlashlightEnablingInnovation2022}. The models tested are a wav2vec2 model finetuned on 10 minutes of audiobook data \citep{panayotovLibrispeechASRCorpus2015}, and Whisper (Table~\ref{tab:asr_reranking}).

\begin{table}[!htbp]
\centering
\caption{WER\% $\downarrow$ of n-best list rescoring strategies including Grad-TTS.}
\begin{tabular}{lccccc}
\toprule
\textbf{AM Model} & {AM Score} & {Grad-TTS LL} & {Oracle Best} & {Oracle worst} & {Random} \\
\midrule 
wav2vec2 & 13.3 & 19.6 & 10.5 & 29.0 & 20.2\\
Whisper & 3.3 & 4.7 & 1.5 & 7.2 & 4.3\\
\bottomrule
\end{tabular}
\label{tab:asr_reranking}
\end{table}
In all cases, the DM likelihoods yield a worse score than the \ac{AM}, although better than the oracle worst and better than random in the wav2vec2 case. This suggests that the likelihoods may have some sensitivity to textual changes but lack an expressive linguistic representation. 
%
To survey how DM likelihood interacts with intelligibility and speaker quality, Grad-TTS is used to adapt TED-LIUM live-talk data \citep{rousseauTEDLIUMAutomaticSpeech2012} to aid an ASR model trained on audiobooks. ASR models perform worse when the acoustic properties differ between train and test data. The spectrogram is passed through a Gaussian blur with kernel 5 and $\sigma=1$, simulating text encoder output $\xblur \approx \enc(\vy)$.\footnote{This method is further explained in Appendix \ref{sec:Grad-TTS encoder and decoder output}} The forward ODE is used to calculate $\x_T$, and $\x_0$ is then produced with the reverse ODE. Given that Grad-TTS is trained on LJSpeech, the newly generated $\x_0$ will be acoustically similar to the data the ASR model was trained on. If Grad-TTS can adapt acoustic-linguistic features, it will reduce WER. Both Euler and RK45 solvers are evaluated (Table~\ref{tab:uda}).

\begin{table}[!htbp]
    \centering
\caption{Quantitative assessment of cross-domain adaptation with Grad-TTS.}
     \begin{tabular}{cccccc}
        \toprule
         \multicolumn{2}{c}{ODE sampler} & & \multicolumn{3}{c}{Metric}\\
         \cmidrule{1-2} \cmidrule{4-6}
         Forward & Reverse & & WER $\downarrow$ & BPD $\downarrow$ & ASV $\uparrow$ \\
         \midrule
          None & None & & 0.46 & 1.66 & 0.52\\
          Euler & Euler & & 0.94 & -1.47 & 0.66\\
          RK45 & Euler & & 0.57 & 6.92 & 0.54\\
          RK45 & RK45 & & 0.57 & 1.83 & 0.55\\
         \bottomrule
    \end{tabular}
    \label{tab:uda}
\end{table}
The Euler method produces unintelligible audio, yet has high speaker similarity and likelihood. The RK45 sampler degrades intelligibility to a lesser degree than Euler, but speaker adaptation is negligible. The trend is that greater shift towards the LJSpeech speaker produces less intelligible audio.

\subsection{Text-to-image}
\label{sec:SDXL Images}

\begin{table}[!htbp]
\centering
\caption{The SDXL generative process from $t=1$ to 0 with CFG=7 at 1024 resolution.}
\label{tab:sdxl_over_t}
\begin{tabular}{l||l|rrrrrrr}
  \toprule
      \multirow{2}*{Metric} & \multirow{2}*{Data} & \multicolumn{5}{c}{Time}\\
                            & & 1.00 & 0.75 & 0.50 & 0.25 & 0.00 \\
  \midrule
      BPD $\downarrow$ & -0.17 & 0.07 & 0.08 & -0.08 & -0.12 & -0.08\\
      FID $\downarrow$ & 0.00 & 478.91 & 369.91 & 225.89 & 221.40 & 229.08\\
      CLIP $\uparrow$ & 29.50 & 22.72 & 23.49 & 28.10 & 27.97 & 28.16\\
  \bottomrule
\end{tabular}
\end{table}
Similar to Table~\ref{tab:gtts_over_t}, sample quality increases~(FID) during the generative process, as does likelihood. Although there is a clear trend that FID does decrease over time, it is higher than typical.\footnote{The DM in \cite{songScoreBasedGenerativeModeling2021} yields 2.92 FID on CIFAR-10.} It is assumed that this is because SDXL was trained on more data than PACS, thus generating many samples absent in the PACS dataset. Unlike Grad-TTS, the conditional metric (CLIP) improves, showing that CFG expresses a more conditional likelihood than the MR-VP SDE used by Grad-TTS.

\begin{table}[!htbp]
  \centering
  \caption{Impacts of CFG, reconstruction (-R) and image sizes.}
  \label{tab:sdxl_cfg_reconstruct}
    \begin{tabular}{l||l|rrrr|rrrr}
      \toprule
          \multirow{2}*{Metric} & \multirow{2}*{Data} & \multicolumn{4}{c|}{CFG} & \multicolumn{4}{c}{Size-Reconstruct}\\
              & & 0.00 & 3.00 & 5.00 & 7.00 & 512 & 512-R & 1024 & 1024-R\\
      \midrule
          BPD $\downarrow$ & -0.17 & -0.04 & -0.05 & -0.07 & -0.08 & 0.02 & -0.01 & 0.07 & -0.01\\
          FID $\downarrow$ & 0.00 &  300.98 & 249.67 & 227.14 & 229.08 & 406.39 & 406.26 & 478.91 & 480.34\\
          CLIP $\uparrow$ & 29.50 & 23.16 & 26.70 & 27.68 & 28.16 & 22.73 & 22.75 & 22.72 & 22.71\\
      \bottomrule
    \end{tabular}
  \end{table}
  Table~\ref{tab:sdxl_cfg_reconstruct} shows that samples produced with higher CFG scale intuitively receive higher likelihood. To measure how the generative process is intertwined with likelihood, the forward process and then the backward process are applied to an image, producing an image similar to the original but with features more consistent with synthetic images. This operation is termed \textit{reconstruction}.\footnote{An example is given in Appendix \ref{sec:SDXL Images}} Reconstruction increases likelihood without affecting other metrics, revealing inconsistency (Table~\ref{tab:sdxl_cfg_reconstruct}).

Lastly, this chapter investigates visual-language reasoning, an important ability for \ac{TTI} models. The aim is to determine whether DM likelihood can correctly match images to captions from a list of compositionally confounding captions. Although diffusion classifiers show that this is possible with proxy-likelihoods, it remains to be determined whether this also holds for exact likelihood. The ability to discern objects is evaluated on CLEVR. For each image, all confounding prompts are presented e.g. the correct prompt is ``a photo of a red sphere'' but instead ``a photo of a red cube'' is provided. Accuracy is calculated based on the number of correct captions that were given the highest likelihood. Both colour and shape are confounded, but not at the same time.
Other visual-language concepts are also considered. Complementary to the CLEVR experiment, PACS is used to measure the likelihood sensitivity to object versus domain changes, for example a photo of a horse.

\begin{table}[!htbp]
  \caption{Prompt accuracy on CLEVR and PACS}\label{tab:clevr_pacs}
  \centering
  \begin{tabular}{rlrrrrrr}
    \toprule
    \multicolumn{2}{c}{Params} && \multicolumn{2}{c}{CLEVR} && \multicolumn{2}{c}{PACS}\\
        \cmidrule{1-2} \cmidrule{4-5} \cmidrule{7-8}
    Reconstruct & CFG && Colour & Shape && Class & Domain\\
    \midrule
    \xmark & 0 && 12 & 33 && 86 & 50 \\
    & 3 &&  0 & 33 && 0 & 0 \\
    & 5 && 25 & 33 && 0 & 0 \\
    % & 7 && 12 & 33 && 0 & 0 \\
    \midrule
    \checkmark & 0 && 38 & 33 && 14 & 0 \\
    & 3 && 12 & 0 && 14 & 0 \\
    & 5 && 12 & 0 && 0  & 0 \\
    % & 7 && 12 & 0 && 0  & 0 \\
    \midrule
    Random &&& 12 & 33 && 11 & 25 \\
    CLIP-score &&& 63 & 67 && 100 & 100 \\
    \bottomrule
  \end{tabular}
\end{table}
The results are shown in Table \ref{tab:clevr_pacs}. It can be seen that the scores are unexpectedly low. Scores from random selection and from selecting captions with the highest CLIP score are also provided. The results on CLEVR are comparable to random choice; when reconstruction is used, shape accuracy is reduced. SDXL performs better on PACS without reconstruction and guidance. This is anomalous behaviour since the rest of the PACS results are typically 0\% accuracy.

\section{Summary}
Given that DMs are trained to maximise the VLB of likelihood, it is important to understand
the nuances of this exact likelihood. This is especially true for conditional tasks where
the conditional likelihood is often modelled implicitly, such as with a MR-VP SDE, or
with CFG. The results show that even reasonable assumptions should be verified, as the task and
conditional ODE used have a substantial impact on the nature of DM likelihood. For Grad-
TTS, the MR-VP SDE likelihood models speaker characteristics and smooth spectrograms,
but not features of language. Such findings are contrary to what one would expect for a
probabilistic conditional process between speech and text. For SDXL, CFG does appear to
model prompt-faithful image generation, but exact likelihood from SDXL cannot be used for
prompt classification. This is inconsistent with the concept of Diffusion Classifiers that use
lower bounds on likelihood for classification. Hence, these quantitative results highlight a current lack
of understanding of DM likelihood. This chapter provides initial evidence that more
attention should be devoted to experiments and theoretical understanding of DM likelihood,
especially on conditional tasks.

\section{Additional tables}

\begin{table}[ht]
  \caption{Various metrics during the generative process over time. The image distribution quality (FID) and prompt accuracy (CLIP) improve, as does the likelihood (BPD). The best results are found at 1024 resolution with a high guidance scale.}
\tiny
    \begin{tabular}{rrrrrrrrrrrrrrrrrrrrr}
\toprule
\multicolumn{3}{c}{pipe} & \multicolumn{3}{c}{1.00} & \multicolumn{3}{c}{0.75} & \multicolumn{3}{c}{0.50} & \multicolumn{3}{c}{0.25} & \multicolumn{3}{c}{0.00} & \multicolumn{3}{c}{ref} \\
CFG & height & r & BPD & FID & CLIP & BPD & FID & CLIP & BPD & FID & CLIP & BPD & FID & CLIP & BPD & FID & CLIP & BPD & FID & CLIP \\
\midrule
0 & 512 & \xmark & 0.02 & 417.58 & 22.58 & 0.02 & 399.30 & 23.62 & -0.01 & 310.19 & 22.68 & -0.01 & 308.96 & 22.23 & -0.01 & 310.24 & 22.55 & -0.04 & 0.00 & 29.56 \\
3 & 512 &  & 0.02 & 412.40 & 22.60 & 0.02 & 391.65 & 23.90 & -0.01 & 280.20 & 23.75 & -0.02 & 279.52 & 23.45 & -0.01 & 278.99 & 23.66 & -0.04 & 0.00 & 29.56 \\
5 & 512 &  & 0.02 & 409.33 & 22.67 & 0.02 & 378.74 & 24.18 & -0.02 & 268.82 & 24.36 & -0.02 & 268.72 & 23.95 & -0.02 & 264.01 & 24.35 & -0.04 & 0.00 & 29.56 \\
7 & 512 &  & 0.02 & 406.39 & 22.73 & 0.02 & 373.26 & 24.31 & -0.02 & 257.55 & 24.63 & -0.03 & 257.93 & 24.31 & -0.02 & 257.00 & 24.56 & -0.04 & 0.00 & 29.56 \\
0 & 512 & \checkmark & 0.01 & 417.80 & 22.58 & 0.00 & 399.60 & 23.61 & -0.01 & 309.89 & 22.65 & -0.02 & 308.74 & 22.23 & -0.01 & 311.61 & 22.57 & -0.04 & 0.00 & 29.56 \\
3 & 512 &  & -0.00 & 412.18 & 22.60 & -0.00 & 392.58 & 23.93 & -0.02 & 282.78 & 23.87 & -0.02 & 278.91 & 23.47 & -0.02 & 280.99 & 23.67 & -0.04 & 0.00 & 29.56 \\
5& 512 &  & -0.01 & 408.25 & 22.68 & -0.01 & 374.78 & 24.19 & -0.02 & 269.79 & 24.41 & -0.03 & 266.02 & 24.08 & -0.02 & 265.90 & 24.28 & -0.04 & 0.00 & 29.56 \\
7 & 512 &  & -0.01 & 406.26 & 22.75 & -0.01 & 373.70 & 24.31 & -0.02 & 256.47 & 24.63 & -0.03 & 256.94 & 24.36 & -0.02 & 259.19 & 24.66 & -0.04 & 0.00 & 29.56 \\
0 & 1024 & \xmark & 0.07 & 478.22 & 22.59 & 0.08 & 418.89 & 21.98 & -0.04 & 300.59 & 23.00 & -0.06 & 309.11 & 22.99 & -0.04 & 300.98 & 23.16 & -0.17 & 0.00 & 29.56 \\
3 & 1024 &  & 0.07 & 478.30 & 22.65 & 0.08 & 399.29 & 22.51 & -0.05 & 245.04 & 26.55 & -0.08 & 251.32 & 26.59 & -0.05 & 249.67 & 26.70 & -0.17 & 0.00 & 29.50 \\
5 & 1024 &  & 0.07 & 478.65 & 22.70 & 0.08 & 380.89 & 23.04 & -0.07 & 231.44 & 27.72 & -0.10 & 232.38 & 27.69 & -0.07 & 227.14 & 27.68 & -0.17 & 0.00 & 29.51 \\
7 & 1024 &  & 0.07 & 478.91 & 22.72 & 0.08 & 369.91 & 23.49 & -0.08 & 225.89 & 28.10 & -0.12 & 221.40 & 27.97 & -0.08 & 229.08 & 28.16 & -0.17 & 0.00 & 29.47 \\
0 & 1024 & \checkmark & 0.06 & 478.22 & 22.59 & 0.08 & 418.82 & 21.97 & -0.04 & 300.53 & 23.02 & -0.06 & 308.12 & 22.99 & -0.04 & 301.76 & 23.17 & -0.17 & 0.00 & 29.56 \\
3 & 1024 &  & 0.05 & 479.91 & 22.65 & 0.03 & 402.99 & 22.49 & -0.06 & 252.33 & 26.55 & -0.09 & 263.05 & 26.60 & -0.06 & 255.12 & 26.77 & -0.17 & 0.00 & 29.48 \\
5 & 1024 &  & 0.02 & 480.88 & 22.68 & 0.00 & 384.40 & 23.01 & -0.07 & 239.35 & 27.69 & -0.11 & 241.89 & 27.66 & -0.07 & 235.62 & 27.73 & -0.18 & 0.00 & 29.47 \\
7 & 1024 &  & -0.01 & 480.34 & 22.71 & -0.01 & 375.69 & 23.43 & -0.09 & 231.81 & 27.96 & -0.13 & 228.23 & 27.95 & -0.09 & 234.73 & 28.10 & -0.18 & 0.00 & 29.48 \\
\bottomrule
\end{tabular}
    \label{tab:sdxl_over_t_full}
\end{table}
